\section{Discussion, Limitations, and Conclusion}
\label{sec:discussion}

% src: cross_model_comparison.json, quantitative_findings_qwen.json, quantitative_findings_gpt54.json
\parbench{} establishes that LLMs can perform non-trivial parallel code translation when evaluation isolates the kernel and preserves executable correctness. Two models show a large capability gap: \gptnew{} achieves pass@1\,=\,62.7\% versus \qwenshort{} at 23.9\% across the same 142 L0 tasks (task-level paired McNemar's $\chi^2 = 45.2$, $p < 10^{-10}$; Cohen's $h = 0.80$; concordance: 49 both-pass, 42 both-fail, 50 GPT-only, 1 Qwen-only). Both models share the same dominant failure mode---incomplete API-surface adaptation at build time---but \gptnew{} reduces build failures from 50.0\% to 22.8\% of L0 records. Direction-level results matter for both models: CUDA-to-OpenMP is the most tractable standard direction (\gptnew{} 83.3\%, \qwenshort{} 40.3\%), while OpenCL-to-CUDA remains the hardest (\gptnew{} 19.3\%, \qwenshort{} 0\%). Augmentation robustness analysis on a balanced 12-kernel subset shows stable pass rates (75--83\%) across five augmentation levels.

\textbf{Limitations.} The benchmark is correctness-focused: it does not measure speedups, occupancy, memory bandwidth, or efficiency. Correct translated code may still be unusably slow, a concern also raised by TRACY~\cite{TRACY2025}. The corpus is finite; some directions have small paired samples, and the augmentation analysis is limited to 12 kernels at L1--L4, precluding formal trend tests. Rodinia and related suites are likely present in training data, and augmentation tests surface-form robustness rather than deeper algorithmic memorization. Ten specs across six kernels were downgraded from strong or medium oracles to weak (stdout pattern plus exit code): six due to cross-API floating-point reduction-order divergence (cfd, hotspot, myocyte) and four due to synthesis asymmetry between reference implementations (bfs, nw, nn). This limits oracle strength but correctly avoids penalizing faithful translations. The stochastic evaluation (temperature~0.7, three samples) captures variability but does not include self-repair; bounded self-repair is a framework capability deferred to future work. The two models use different reasoning mechanisms and sampling configurations (Section~\ref{sec:sampling-config}). The observed gap is large (Cohen's $h = 0.80$) but cannot be fully attributed to model capability given unmatched sampling conditions; it reflects an unknown mixture of model strength and provider-side differences.

\textbf{Future work.} Additional directions include self-repair experiments using the framework's built-in feedback loop, evaluation of agentic systems such as LASSI~\cite{LASSI2024}, fine-tuned HPC models~\cite{HPCCoderV2}, efficiency-aware evaluation following TRACY~\cite{TRACY2025}, and expansion of the augmentation campaign to more kernels and directions.

\parbench{} is evaluation infrastructure rather than a new translation method. Its contribution is to make the measurement problem sharper: kernel-centric isolation avoids repository-level build confounds, declarative specs make correctness reproducible, augmentation probes robustness, and the harness exposes where translations fail. The initial results show that current LLMs have real but uneven parallel translation capability, with build adaptation, direction asymmetry, and multi-file coordination as central barriers to reliable use.
